\worldchapter{Magie}{magic}
\label{ch:magie}

\begin{multicols*}{2}

\section{Magisches Tirakan}

Tirakan ist eine magische Welt.
Jedes Lebewesen auf Tirakan hat eine gewisse Fähigkeit, mit der Magie der Welt umzugehen.
Ein Charakter kann die magische Kraft kontrollieren und Zauber sprechen, wenn er mindestens einen Punkt in der Fähigkeit \textbf{Zauberei} hat.

Für einen Charakter gilt:

\begin{itemize}[leftmargin=*]
\item Zauberei >= 1: eine Schule wählen
\item Zauberei >= 2: zwei Schulen wählen
\item Startzauber: 2 + Zauberei Rang, maximal 5
\item bekannte Zauber müssen aus den gewählten Schulen stammen
\end{itemize}

Schulen der Magie sind:

\begin{itemize}[leftmargin=*]
\item Astrale Magie
\item Chimärologie
\item Dämonologie
\item Elementarmagie
\item Geisterbeschwörungen
\item Hermetik
\item Nekrologie
\item Mystizismus
\item Sanguine Magie
\item Schamanismus
\item Schwarze Magie
\item Verfluchungen
\item Weiße Magie
\item Zauberei
\end{itemize}

\subsection{Wirkvorgang}
\begin{enumerate}[leftmargin=*]
\item einen bekannten Zauber wählen
\item Wirkmodus wählen
\item Stufe wählen
\item Mind oder Will + Zauberei Probe
\item Arcana-Pool-Kosten immer bezahlen, auch bei Fehlschlag
\end{enumerate}

\subsection{Wirkmodi}
\begin{itemize}[leftmargin=*]
\item \textbf{Fast Cast}\index{Fast Cast}: 1 Action
\item \textbf{Ritual Cast}\index{Ritual Cast}: 2+ Aktionen oder ein Auszeit-Segment, wenn Vorbereitung, Kreise, Hilfen oder Material nötig sind
\end{itemize}

\subsection{Wirkstufen}
\begin{description}[style=nextline,leftmargin=0pt]
\item[\textbf{Minor}]\index{Minor} +10 Zielwert, kein Zuschlag
\item[\textbf{Severe}]\index{Severe} -10 Zielwert, +1 Arcana
\item[\textbf{Catastrophic}]\index{Catastrophic} -30 Zielwert, +2 Arcana
\end{description}

\subsection{Arcana-Kosten}
Basiskosten nach spell\_points:
\begin{itemize}[leftmargin=*]
\item 1--3 \textrightarrow{} 1 Arcana
\item 4--7 \textrightarrow{} 2 Arcana
\item 8--11 \textrightarrow{} 3 Arcana
\item 12+ \textrightarrow{} 4 Arcana
\end{itemize}

Gesamtkosten sind Basis plus Stufenzuschlag, Minimum 1.

\subsection{Ergebnisse}
\begin{itemize}[leftmargin=*]
\item Kritischer, starker oder sauberer Erfolg: Effekt tritt wie erklärt ein
\item Knapp geschafft: Effekt tritt ein, aber mit einem kleinen instabilen Detail
\item jeder Misserfolg: kein beabsichtigter Effekt, stattdessen Backlash
\end{itemize}

\subsection{Zauberverstärkung}
Bei erfolgreicher Zauberei gibt \textbf{Magic Level}\index{Magic Level} genau \textbf{einen} Boost-Schritt auf:
\begin{itemize}[leftmargin=*]
\item Reichweite
\item Dauer
\item Zielzahl
\item Schaden- oder Clear-Wert
\end{itemize}

\subsection{Backlash}\index{Backlash}
\begin{itemize}[leftmargin=*]
\item \textbf{Minor}: +1 Burden oder \textbf{Shaken}
\item \textbf{Serious}: +1 Wound, großer Kollateralschaden oder feindliche Aufmerksamkeit
\item \textbf{Critical}: +2 Burden, schwerer Kollateralschaden, Riss oder bleibender Fluchmal-Effekt
\end{itemize}

Bei einem misslungenen 96--00 kommt zusätzlich ein thematischer Nebeneffekt hinzu.

\subsection{Chaotische Zauberei}
Chaotische Zauberei gilt, wenn:
\begin{itemize}[leftmargin=*]
\item Magic Level >= 5, oder
\item der zaubernde Charakter 2+ Verderbnis hat
\end{itemize}

Nach jedem Wirken 1d6:
\begin{enumerate}[leftmargin=*]
\item Echo: Effekt trifft ein nächstes gültiges Nebenziel mit Limited Effect
\item Drain: 1 zusätzliches Arcana verlieren
\item Fracture: nahe Zerbrechlichkeit oder Ward bricht
\item Beacon: okkulte Aufmerksamkeit richtet sich auf die Zaubernden
\item Reversal Pulse: ein kleines beabsichtigtes Detail kehrt sich um
\item Stable Surge: kein Zusatzeffekt
\end{enumerate}

\end{multicols*}
